% ============================================================================
%  Appendice C -- Riferimento di Git e GitHub
% ============================================================================
\chapter{Riferimento di Git e GitHub}
\label{app:github}

La piattaforma di questo libro, gli strumenti, il manuale di accompagnamento e
il libro stesso vivono su GitHub, e lì dovrebbe vivere anche il vostro
progetto. Questa appendice è un riferimento di Git, il programma che registra
la storia di un progetto, e di GitHub, il servizio che lo ospita: i concetti,
i comandi di ogni giorno, il flusso di lavoro con cui si contribuisce a \prg,
e il modo in cui si conserva e si rilascia un progetto come \vermone. Non
presuppone alcuna esperienza. Le fonti autorevoli sono \citet{progit} e la
documentazione di GitHub \citep{githubdocs}.

% ----------------------------------------------------------------------------
\section{Perché il controllo di versione}

Il Capitolo~\ref{ch:verm1p1} ha dato un'istruzione che ogni tutorial ha
ripetuto: \emph{tenete una copia di ogni parte finita}. Una cartella piena di
file chiamati \fname{verm1\_part2\_ok.S} e \fname{verm1\_final\_really.S} è la
forma artigianale del controllo di versione. Git fa la stessa cosa come si
deve: ricorda ogni stato registrato dei file, dice esattamente che cosa è
cambiato tra due stati, e ritorna a uno qualsiasi di essi.

\begin{prgconcept}
\textbf{Git} registra istantanee di un progetto, chiamate \term{commit};
ciascuna è uno stato etichettato al quale si può tornare. \textbf{GitHub}
ospita in linea un repository Git, così che il lavoro ha una copia di
sicurezza, può essere condiviso e può essere rivisto. Insieme trasformano
«credo di aver rotto qualcosa» in «fammi vedere esattamente che cosa ho
cambiato».
\end{prgconcept}

% ----------------------------------------------------------------------------
\section{Concetti}

\begin{table}[htb]
\centering\small
\caption{Il vocabolario di Git e GitHub.}\label{tab:git:terms}
\begin{tabularx}{\textwidth}{@{}lL@{}}
\toprule
\textbf{Termine} & \textbf{Significato}\\
\midrule
repository & un progetto con tutta la sua storia: la contiene la cartella \fname{.git} dentro il progetto\\
working tree & i file come sono ora nella cartella\\
staging area (index) & le modifiche scelte per il prossimo commit\\
commit & un'istantanea registrata, con un autore, una data, un messaggio e un identificatore unico (un hash di 40 cifre esadecimali, di solito abbreviato a 7)\\
branch & un nome mobile per una linea di commit; quello predefinito è \code{main}\\
\code{HEAD} & il commit su cui è basato il working tree\\
tag & un nome fisso per un commit, tipicamente una versione: \code{v1.0.0}\\
remote & una copia del repository altrove, tipicamente su GitHub; il nome predefinito è \code{origin}\\
clone & una copia locale completa di un repository remoto\\
fork & una copia del repository di qualcun altro sotto il proprio account GitHub\\
push / fetch / pull & inviare commit a un remote / riceverli / riceverli e fonderli\\
merge & unire le storie di due branch\\
pull request (PR) & una proposta, su GitHub, di fondere un branch in un altro, con una discussione e una revisione\\
issue & una segnalazione o un compito associato a un repository su GitHub\\
release & un tag pubblicato su GitHub con note e, opzionalmente, file da scaricare\\
\bottomrule
\end{tabularx}
\end{table}

\begin{figure}[htb]
\centering
\begin{tikzpicture}[node distance=7mm]
  \node[blkacc,text width=27mm,minimum height=14mm] (w) {\textbf{working tree}\\\footnotesize i file che modificate};
  \node[blkamb,text width=27mm,minimum height=14mm,right=of w] (s) {\textbf{staging area}\\\footnotesize il prossimo commit};
  \node[blkvio,text width=27mm,minimum height=14mm,right=of s] (l) {\textbf{repository locale}\\\footnotesize la storia};
  \node[blkgrn,text width=27mm,minimum height=14mm,right=of l] (r) {\textbf{remote}\\\footnotesize GitHub};
  \draw[->,thick] (w) -- node[above,font=\ttfamily\scriptsize]{add} (s);
  \draw[->,thick] (s) -- node[above,font=\ttfamily\scriptsize]{commit} (l);
  \draw[->,thick] ([yshift=2mm]l.east) -- node[above,font=\ttfamily\scriptsize]{push} ([yshift=2mm]r.west);
  \draw[->,thick] ([yshift=-2mm]r.west) -- node[below,font=\ttfamily\scriptsize]{fetch} ([yshift=-2mm]l.east);
  \draw[->,thick,rounded corners=3mm] (l.south) -- ++(0,-7mm) -| node[pos=0.25,below,font=\ttfamily\scriptsize]{restore, switch} (w.south);
\end{tikzpicture}
\caption{I quattro luoghi in cui può trovarsi una modifica, e i comandi che la spostano.}
\label{fig:git:areas}
\end{figure}

Una modifica viaggia attraverso quattro luoghi (Figura~\ref{fig:git:areas}).
Viene fatta nel working tree; \code{git add} la sceglie per il prossimo
commit; \code{git commit} la registra nel repository locale; \code{git push}
la invia al remote. Niente lascia il computer prima del push.

% ----------------------------------------------------------------------------
\section{Installare e configurare}

Git si installa con il gestore di pacchetti del sistema
(\code{brew install git} su macOS, \code{sudo apt install git} su Debian e
Ubuntu) oppure, su Windows, con \emph{Git for Windows} da
\url{https://git-scm.com}; è già stato installato nel
Capitolo~\ref{ch:setup}. Si configura una volta per computer
(Listato~\ref{lst:git:config}).

\begin{shell}[caption={Configurare Git, una volta.},label=lst:git:config]
git --version
git config --global user.name  "Your Name"
git config --global user.email "your.address@studenti.uniparthenope.it"
git config --global init.defaultBranch main
git config --list
\end{shell}

Un account GitHub si crea su \url{https://github.com}. Per far parlare Git con
GitHub senza scrivere una password a ogni push, il modo più semplice è il
client ufficiale a riga di comando, \code{gh}, il cui comando
\code{gh auth login} apre un browser e memorizza la credenziale nel portachiavi
del sistema; l'alternativa è una chiave SSH (\code{ssh-keygen -t ed25519}), la
cui metà \emph{pubblica} si incolla nelle impostazioni dell'account.

\begin{prgwarn}
Una chiave privata, una password e un token non si mettono mai in un commit,
non si incollano in un file del progetto e non si mostrano in una schermata.
Questo vale per il token del Cartridge Store esattamente come per le
credenziali di GitHub.
\end{prgwarn}

% ----------------------------------------------------------------------------
\section{Ottenere un repository}

Ci sono due modi di ottenere un repository: copiarne uno esistente, o
trasformare una cartella in un repository (Listato~\ref{lst:git:get}).

\begin{shell}[caption={Clonare la piattaforma; avviare un proprio progetto.},label=lst:git:get]
# 1. clone an existing repository: the platform, the debugger, the manual
git clone https://github.com/riscv-prg32/PRG32
git clone https://github.com/riscv-prg32/PRG32-QT
git clone https://github.com/riscv-prg32/Playing-with-RISC-V_PRG32_textbook

# 2. start a new repository in the folder of your game
cd ~/verm1
git init
git add verm1.S
git commit -m "Verm1 part 1: a head that crosses the world"
\end{shell}

Per aggiornare un clone rispetto al remote: \code{git pull}.

% ----------------------------------------------------------------------------
\section{Il ciclo di ogni giorno}

Tre comandi coprono la maggior parte dell'uso quotidiano: \code{status}, per
vedere che cosa è cambiato; \code{add}, per scegliere che cosa tenere;
\code{commit}, per registrarlo con un messaggio
(Listato~\ref{lst:git:cycle}).

\begin{shell}[caption={Ispezionare, preparare, registrare.},label=lst:git:cycle]
git status                      # what has changed?
git diff                        # exactly which lines?
git add verm1.S                 # choose the changes to record
git diff --staged               # what is about to be recorded?
git commit -m "Part 2: steering, with the U-turn rule as one XOR"
git log --oneline               # the history, one line per commit
\end{shell}

\begin{prgconcept}[Un commit dopo ogni \emph{Eseguite}]
Il metodo di questo libro è un passo eseguibile alla volta. Fate un commit a
ogni punto di controllo in cui il programma funziona e si comporta come
dichiarato: allora ogni rottura è, nel peggiore dei casi, il lavoro di un
passo da ispezionare o da annullare.
\end{prgconcept}

Un messaggio di commit è una nota per sé stessi nel futuro. La sua prima riga
dice all'imperativo che cosa fa il commit, in una cinquantina di caratteri:
«Add the edge detector for button A», non «stuff» o «fix». Se serve altro,
segue una riga vuota e poi una spiegazione del \emph{perché}.

\subsection{Che cosa non registrare: \texttt{.gitignore}}

I file generati non appartengono a un repository: possono essere ricostruiti
dai sorgenti, e cambiano a ogni compilazione. Un file di nome
\fname{.gitignore} elenca i modelli che Git deve ignorare
(Listato~\ref{lst:git:ignore}).

\begin{plain}[caption={Un \fname{.gitignore} per un progetto di cartuccia.},label=lst:git:ignore]
# build products
build/
*.o
*.elf
*.prg32
*.zip

# editors and operating systems
.vscode/
.DS_Store
Thumbs.db

# secrets: never
.env
*.token
\end{plain}

Se tenere sotto controllo di versione la cartuccia compilata è una questione
di scelta; la risposta usuale è no, perché il posto giusto per essa è una
release (Paragrafo~\ref{sec:git:release}).

% ----------------------------------------------------------------------------
\section{Guardare la storia}

\begin{shell}[caption={Leggere la storia.},label=lst:git:log]
git log --oneline --graph --all     # all the branches, as a graph
git log -p verm1.S                  # the changes made to one file
git show 3f2a9c1                    # one commit: message and changes
git diff v1.0.0 HEAD                # from a version to now
git blame verm1.S                   # who wrote each line, and when
git bisect start                    # find the commit that broke something
\end{shell}

\code{git bisect} merita una parola. Dato un commit in cui il programma
funzionava e uno in cui non funziona, propone il commit a metà strada, chiede
se quello è buono o cattivo, e ripete: una ricerca binaria nella storia, che
trova il commit colpevole tra mille in dieci passi.

% ----------------------------------------------------------------------------
\section{Annullare}

\begin{table}[htb]
\centering\small
\caption{Come annullare, dal meno al più drastico.}\label{tab:git:undo}
\begin{tabularx}{\textwidth}{@{}LP{0.44\textwidth}@{}}
\toprule
\textbf{Situazione} & \textbf{Comando}\\
\midrule
scartare le modifiche di un file dall'ultimo commit & \code{git restore verm1.S}\\
togliere un file dalla staging area, conservandone le modifiche & \code{git restore -{}-staged verm1.S}\\
correggere il messaggio, o aggiungere un file dimenticato, dell'ultimo commit non ancora inviato & \code{git commit -{}-amend}\\
riportare un file a com'era in un commit più vecchio & \code{git restore -{}-source 3f2a9c1 verm1.S}\\
cancellare l'effetto di un commit già pubblicato, con un nuovo commit & \code{git revert 3f2a9c1}\\
mettere da parte il lavoro per un momento, e riprenderlo & \code{git stash} ; \code{git stash pop}\\
guardare un vecchio stato senza cambiare niente & \code{git switch -{}-detach 3f2a9c1}\\
\bottomrule
\end{tabularx}
\end{table}

\begin{prgwarn}
I comandi che riscrivono la storia, come \code{git reset -{}-hard} e
\code{git push -{}-force}, distruggono lavoro, il proprio e quello degli altri.
La storia che è stata inviata si corregge con \code{git revert}, che aggiunge
un commit, e mai riscrivendola.
\end{prgwarn}

% ----------------------------------------------------------------------------
\section{Branch}

Un \term{branch} (ramo) è una linea di lavoro indipendente. Prima di una
modifica rischiosa, per esempio il disegno incrementale del
Capitolo~\ref{ch:verm1p4}, si crea un branch: se l'esperimento fallisce il
branch si scarta e la linea principale resta intatta; se riesce viene fuso
(Listato~\ref{lst:git:branch}).

\begin{shell}[caption={Lavorare su un branch.},label=lst:git:branch]
git switch -c feature/pause-state   # create a branch and move to it
# ... edit, build, run, commit, as many times as needed ...
git switch main                     # back to the main line
git merge feature/pause-state       # bring the work in
git branch -d feature/pause-state   # the branch is no longer needed

git branch                          # list the branches
git switch -                        # go back to the previous branch
\end{shell}

\paragraph{Conflitti.} Se due branch hanno modificato le stesse righe, il
merge si ferma e Git segna il punto nel file:
\begin{plain}
<<<<<<< HEAD
.equ SPEED_MIN, 32
=======
.equ SPEED_MIN, 40
>>>>>>> feature/levels
\end{plain}
Il conflitto si risolve modificando il file finché dice ciò che si vuole,
togliendo le tre righe di marcatura, e poi eseguendo \code{git add} e
\code{git commit}. \code{git merge -{}-abort} rinuncia e ritorna allo stato
precedente al merge.

% ----------------------------------------------------------------------------
\section{Lavorare con GitHub}

\subsection{Pubblicare un repository}

\begin{shell}[caption={Pubblicare un repository locale su GitHub.},label=lst:git:publish]
# with the GitHub client: create the repository and push, in one step
gh repo create verm1-yourname --public --source . --push

# or: create an empty repository on github.com, then
git remote add origin https://github.com/yourname/verm1-yourname.git
git push -u origin main

# afterwards
git push            # send new commits
git pull            # receive the commits made elsewhere
git remote -v       # where does this repository point?
\end{shell}

\subsection{I file di un buon repository}

Il repository di \prg è un modello delle convenzioni
(Tabella~\ref{tab:git:files}): i file che orientano chi arriva e rendono un
progetto facile da valutare.

\begin{table}[htb]
\centering\small
\caption{I file che un repository dovrebbe contenere.}\label{tab:git:files}
\begin{tabularx}{\textwidth}{@{}lL@{}}
\toprule
\textbf{File} & \textbf{Contenuto}\\
\midrule
\fname{README.md} & che cos'è il progetto, una schermata, come compilarlo ed eseguirlo, i comandi\\
\fname{LICENSE} & che cosa possono farne gli altri: MIT per \prg e per i programmi di questo libro\\
\fname{.gitignore} & che cosa non va registrato\\
\fname{CITATION.cff} & come citare il progetto in una relazione o in una tesi\\
\fname{CONTRIBUTING.md} & come proporre una modifica\\
\fname{CHANGELOG.md} & che cosa è cambiato da una versione alla successiva\\
\fname{.github/workflows/} & i controlli automatici eseguiti da GitHub Actions\\
\bottomrule
\end{tabularx}
\end{table}

\begin{plain}[caption={Un \fname{README.md} per una variante di \vermone.},label=lst:git:readme]
# Verm1 Deluxe

A hungry worm on a 40 by 24 grid: a tribute to Snake, written in RISC-V
assembly for the PRG32 platform. My variant adds walls and a pause.

![The game](screenshot.png)

## Build and run

    python3 -m prg32 cartridge build verm1.S --portable \
        --entry-prefix verm1 --name verm1 --out verm1.prg32
    python3 -m prg32 esp32c6 upload-and-run verm1.prg32 \
        --url http://localhost:8080 --slot cart0

## Controls

Joystick: steer. A: start. B: pause.

## Credits and licence

Written by Your Name for ARC12, University of Naples "Parthenope".
MIT License: see LICENSE.
\end{plain}

I file con l'estensione \fname{.md} sono scritti in Markdown: una riga che
comincia con \code{\#} è un titolo, \code{**bold**} è grassetto, una riga
rientrata di quattro spazi è codice, e \code{[testo](indirizzo)} è un
collegamento.

\subsection{Fork e pull request}

Non si può fare push su un repository di cui non si è manutentori. Il modo di
proporre una modifica al progetto di qualcun altro è la \term{pull request}
(Figura~\ref{fig:git:pr}, Listato~\ref{lst:git:pr}).

\begin{figure}[htb]
\centering
\begin{tikzpicture}[node distance=4mm]
  \node[blkacc,text width=21mm,minimum height=15mm] (a) {\textbf{Fork}\\[1pt]\footnotesize la vostra copia su GitHub};
  \node[blkamb,text width=21mm,minimum height=15mm,right=of a] (b) {\textbf{Branch}\\[1pt]\footnotesize uno scopo chiaro};
  \node[blkvio,text width=21mm,minimum height=15mm,right=of b] (c) {\textbf{Commit, push}\\[1pt]\footnotesize sul vostro fork};
  \node[blkgrn,text width=21mm,minimum height=15mm,right=of c] (d) {\textbf{Pull request}\\[1pt]\footnotesize proporre, discutere};
  \node[blkred,text width=21mm,minimum height=15mm,right=of d] (e) {\textbf{Revisione, merge}\\[1pt]\footnotesize di un manutentore};
  \foreach \x/\y in {a/b,b/c,c/d,d/e} \draw[->,thick] (\x) -- (\y);
\end{tikzpicture}
\caption{Il percorso di un contributo.}
\label{fig:git:pr}
\end{figure}

\begin{shell}[caption={Un contributo a \prg, dal fork alla pull request.},label=lst:git:pr]
gh repo fork riscv-prg32/PRG32 --clone       # fork, and clone the fork
cd PRG32
git switch -c docs/cartridge-workflow-notes  # a branch from main

# ... make the smallest change that proves the idea; run the checks ...

git add docs/
git commit -m "Clarify the cartridge workflow for first-time users"
git push -u origin docs/cartridge-workflow-notes
gh pr create --fill                          # open the pull request

# keep the fork up to date with the original
git remote add upstream https://github.com/riscv-prg32/PRG32.git
git fetch upstream
git switch main
git merge upstream/main
\end{shell}

Le regole del repository di \prg, nel suo \fname{CONTRIBUTING.md}, sono
scritte prima di tutto per gli studenti: scegliere uno scopo chiaro; creare un
branch da \code{main}; fare la modifica più piccola che dimostri l'idea;
aggiornare la documentazione quando cambiano comportamento, comandi, percorsi o
interfacce; eseguire i controlli; e aprire una pull request che dica qual era
il \emph{contesto}, quale \emph{metodo} è stato seguito e come la modifica è
stata \emph{validata}, su QEMU o sull'hardware. I nomi dei branch dicono il
tipo di lavoro: \code{lab/\ldots}, \code{game/\ldots}, \code{docs/\ldots}

\paragraph{La revisione.} Una pull request mostra le modifiche riga per riga,
e chiunque può commentare una riga. Un revisore approva, o chiede modifiche;
l'autore risponde inviando altri commit sullo stesso branch, che aggiornano la
pull request. La lista di controllo del Capitolo~\ref{ch:store}, usata da un
collega sul vostro bundle, è una revisione.

\subsection{Issue}

Una \term{issue} segnala un difetto o propone un compito. Una buona
segnalazione di un difetto contiene ciò che è stato fatto (il comando esatto),
ciò che ci si aspettava, ciò che è successo (l'\emph{intero} messaggio, come
testo), e l'ambiente: il sistema operativo, la versione degli strumenti,
l'host. Un commit o una pull request il cui messaggio contiene
\code{Fixes \#12} chiude la issue 12 quando viene fuso.

\subsection{GitHub Actions: controlli che girano da soli}

Un file in \fname{.github/workflows} descrive un lavoro che GitHub esegue a
ogni push o pull request, su una macchina pulita: tipicamente, compilare il
progetto ed eseguirne i test. Il workflow del repository di \prg compila le
cartucce incluse con lo stesso strumento a riga di comando che si usa in
locale, per l'ESP32-C6 e per QEMU, e le espone come file da scaricare. Un
segno verde accanto a un commit significa che i controlli sono passati; una
croce rossa significa che qualcosa che funzionava sul computer dell'autore non
funziona su uno pulito, che è esattamente ciò che si vuole sapere prima di un
merge.

\begin{plain}[caption={Lo schema di un file di workflow (YAML). Il passo di compilazione
dipende dal progetto: per una cartuccia richiede la toolchain ESP-IDF.},label=lst:git:actions]
name: build
on: [push, pull_request]
jobs:
  build:
    runs-on: ubuntu-latest
    steps:
      - uses: actions/checkout@v4
      - name: Build and test
        run: ./scripts/build-and-test.sh
\end{plain}

% ----------------------------------------------------------------------------
\section{Versioni, tag e release}\label{sec:git:release}

Il Capitolo~\ref{ch:store} ha dato a \vermone una versione, 1.0.0, con la
regola che una versione pubblicata non viene mai sostituita. La numerazione si
chiama \term{versionamento semantico}: \emph{major.minor.patch}, dove il patch
cresce per una correzione, il minor per una nuova funzione che non rompe
niente, e il major per una modifica non compatibile. In Git una versione è un
\term{tag}, e su GitHub un tag diventa una \term{release}, alla quale si
possono allegare file (Listato~\ref{lst:git:release}).

\begin{shell}[caption={Etichettare la versione 1.0.0 e rilasciare il bundle.},label=lst:git:release]
git tag -a v1.0.0 -m "Verm1 1.0.0: first release"
git push origin v1.0.0

# publish a release on GitHub and attach the bundle of the Store
gh release create v1.0.0 verm1-1.0.0.zip \
    --title "Verm1 1.0.0" --notes "First release: the finished game."

git tag                 # list the versions
git switch --detach v1.0.0   # look at the sources of a version
\end{shell}

Il tag garantisce che i sorgenti da cui è stato costruito un bundle pubblicato
si possano sempre ritrovare: lega il file che i giocatori installano alle
righe che l'autore ha scritto.

% ----------------------------------------------------------------------------
\section{Questo libro e Overleaf}

I sorgenti di questo libro sono una cartella di file \LaTeX{} che si compila
su Overleaf (caricate la cartella come file ZIP, o importate il repository; il
documento principale è \fname{main.tex} e il compilatore è pdf\LaTeX).
Overleaf può essere collegato a un repository GitHub, con il quale scambia
commit; e ogni progetto Overleaf è esso stesso un repository Git che si può
clonare, così che il testo può essere modificato in un browser e versionato
allo stesso modo di un programma. Le slide del corso sono generate da sorgenti
conservati alla stessa maniera, e le correzioni sono benvenute come pull
request.

% ----------------------------------------------------------------------------
\section{Citare il software}

Il software usato in una relazione, in una tesi o in un articolo si cita, come
qualsiasi altra fonte. Un repository che contiene un file
\fname{CITATION.cff} mostra su GitHub un pulsante, \emph{Cite this
repository}, che produce il riferimento nei formati usuali. Il repository di
\prg ne ha uno, e la sua voce nella bibliografia di questo libro
\citep{prg32repo} è il modello da seguire.

% ----------------------------------------------------------------------------
\section{Riepilogo dei comandi}

{\small
\begin{longtable}{@{}p{0.43\textwidth}p{0.51\textwidth}@{}}
\caption{I comandi di Git e del client GitHub usati in questa appendice.}\label{tab:git:summary}\\
\toprule
\textbf{Comando} & \textbf{Effetto}\\
\midrule
\endfirsthead
\toprule
\textbf{Comando} & \textbf{Effetto}\\
\midrule
\endhead
\bottomrule
\endfoot
\code{git config -{}-global user.name "\ldots"} & impostare il nome dell'autore\\
\code{git init} & fare della cartella corrente un repository\\
\code{git clone URL} & copiare un repository remoto\\
\code{git status} & che cosa è cambiato\\
\code{git diff} ; \code{git diff -{}-staged} & le modifiche non ancora preparate; quelle preparate\\
\code{git add FILE} ; \code{git add -p} & preparare un file; preparare una parte delle modifiche\\
\code{git commit -m "message"} & registrare un'istantanea\\
\code{git log -{}-oneline -{}-graph} & la storia\\
\code{git show COMMIT} & un commit\\
\code{git restore FILE} & scartare le modifiche di un file\\
\code{git restore -{}-staged FILE} & togliere un file dalla staging area\\
\code{git revert COMMIT} & annullare un commit con un nuovo commit\\
\code{git stash} ; \code{git stash pop} & mettere da parte il lavoro; riprenderlo\\
\code{git switch -c NAME} & creare un branch e spostarvisi\\
\code{git switch NAME} & spostarsi su un branch\\
\code{git merge NAME} & fondere un branch in quello corrente\\
\code{git branch} ; \code{git branch -d NAME} & elencare i branch; cancellarne uno\\
\code{git remote -v} ; \code{git remote add NAME URL} & elencare i remote; aggiungerne uno\\
\code{git fetch} ; \code{git pull} & ricevere commit; ricevere e fondere\\
\code{git push} ; \code{git push -u origin NAME} & inviare commit; il primo push di un branch\\
\code{git tag -a v1.0.0 -m "\ldots"} & dare un nome a una versione\\
\code{git bisect start} / \code{good} / \code{bad} & trovare il commit che ha introdotto un difetto\\
\midrule
\code{gh auth login} & accedere a GitHub dal terminale\\
\code{gh repo create NAME -{}-source . -{}-push} & creare un repository su GitHub dalla cartella corrente\\
\code{gh repo fork OWNER/REPO -{}-clone} & fare il fork di un repository e clonare il fork\\
\code{gh pr create} ; \code{gh pr status} ; \code{gh pr checkout N} & aprire una pull request; vederne lo stato; provarne una in locale\\
\code{gh issue create} ; \code{gh issue list} & aprire una issue; elencarle\\
\code{gh release create TAG FILES} & pubblicare una release\\
\end{longtable}}

\begin{prgcheck}
Ora dovreste saper clonare la piattaforma, registrare il vostro lavoro passo
per passo con messaggi significativi, sperimentare su un branch, annullare una
modifica indesiderata, pubblicare il vostro repository, proporre una modifica
con una pull request e rilasciare una versione del vostro gioco. Queste
operazioni bastano per un primo anno di corso, e per un primo contributo a un
progetto aperto.
\end{prgcheck}
